% GENERATED FILE. Edit canonical lessons, then run npm run generate:books.
% canonical-source-hash: fnv1a64:9b59fa48c95c36f5
% canonical-lessons: PA-C216-gae-return, PA-C216-gia-gai, PA-C216-gia-review, PA-C216-kall, PA-C216-where-did-they-go, PA-C216-payoff

\chapter{Where Did They Go?}
\label{ch:pa-a2-gia}
\hlchaptermodality{216}

\begin{chapteropening}
\textbf{By the end of this chapter:} \emph{I can say where someone went, yesterday or earlier, with gia, gai or gae agreeing with them.}

Say who went where, with gia for a man, gai for a woman and gae for a group, and use kall for yesterday.
\end{chapteropening}

\section[\texorpdfstring{\pa{ਉਹ} \pa{ਗਏ} (uh gae)}{ਉਹ ਗਏ (uh gae)}]{Return to a familiar went}

\label{lesson:PA-C216-gae-return}



\begin{quote}
\emph{Hear:} \emph{oh gae} Say \textbf{\pa{ਉਹ} \pa{ਗਏ}}, \textquotedblleft{}they went\textquotedblright{}. Then \emph{to go}: \textbf{\pa{ਜਾਣਾ}}.
\end{quote}

\begin{grammarlens}[title={Grammar Lens}]
\textbf{\pa{ਗਏ}} is the went form for a group, or for one person you respect. Today it gets two partners: one for a man, one for a woman. Hold \textbf{\pa{ਉਹ} \pa{ਰੋਜ਼} \pa{ਸਕੂਲ} \pa{ਜਾਂਦਾ} \pa{ਹੈ}} beside it: every day, and once, in the past.
\end{grammarlens}

\subsection*{Your turn}
\begin{itemize}
\raggedright
  \item \emph{Hear:} \emph{they went} Say \textbf{\pa{ਉਹ} \pa{ਗਏ}}.
  \item \emph{Hear:} \emph{he goes to school every day} Say \textbf{\pa{ਉਹ} \pa{ਰੋਜ਼} \pa{ਸਕੂਲ} \pa{ਜਾਂਦਾ} \pa{ਹੈ}}.
\end{itemize}

\begin{tcolorbox}[breakable,colback=teal!4,colframe=teal!35!black,title={Before you move on}]
Say \emph{they went} once more. No new form is needed yet.
\end{tcolorbox}
\section[\texorpdfstring{\pa{ਗਿਆ} · \pa{ਗਈ} (giā · gaī)}{ਗਿਆ · ਗਈ (giā · gaī)}]{\pa{ਗਿਆ}, \pa{ਗਈ} — went}

\label{lesson:PA-C216-gia-gai}



\begin{quote}
Say \emph{they went}. Then \emph{home} and \emph{school}.
\end{quote}

\begin{grammarlens}[title={Grammar Lens: one went for each}]
The went form agrees with whoever went: \textbf{\pa{ਗਿਆ}} for a man or a boy, \textbf{\pa{ਗਈ}} for a woman or a girl, and the \textbf{\pa{ਗਏ}} you know for a group.

\begin{quote}
\textbf{\pa{ਅਮਨ} \pa{ਘਰ} \pa{ਗਿਆ।}} — \emph{Aman ghar giā.} — Aman went home.
\end{quote}

\begin{quote}
\textbf{\pa{ਕੁੜੀ} \pa{ਸਕੂਲ} \pa{ਗਈ।}} — \emph{kuṛī sakūl gaī.} — The girl went to school.
\end{quote}

\begin{quote}
\textbf{\pa{ਉਹ} \pa{ਕਿੱਥੇ} \pa{ਗਏ}?} — \emph{uh kitthe gae?} — Where did they go?
\end{quote}

There is no helper at the end. CDIAL lists Punjabi \emph{geā}, \textquotedblleft{}gone\textquotedblright{}, under Sanskrit \emph{gata}, \textquotedblleft{}gone\textquotedblright{} (CDIAL 4008).
\end{grammarlens}

\subsection*{Your turn}
\begin{itemize}
\raggedright
  \item \emph{Hear:} \emph{the boy went to school} Say \textbf{\pa{ਮੁੰਡਾ} \pa{ਸਕੂਲ} \pa{ਗਿਆ।}}
  \item \emph{Hear:} \emph{she went home} Say \textbf{\pa{ਉਹ} \pa{ਘਰ} \pa{ਗਈ।}}
  \item \emph{Read it:} \textbf{\pa{ਅਮਨ} \pa{ਘਰ} \pa{ਗਿਆ।}} Give its meaning.
  \item \emph{Write it:} copy \textbf{\pa{ਗਿਆ}} and \textbf{\pa{ਗਈ}} once
\end{itemize}

\begin{tcolorbox}[breakable,colback=teal!4,colframe=teal!35!black,title={Before you move on}]
\emph{Went}, for a man? (\textbf{\pa{ਗਿਆ}}.) For a woman? (\textbf{\pa{ਗਈ}}.) For a group? (\textbf{\pa{ਗਏ}}.)
\end{tcolorbox}
\section[\texorpdfstring{\pa{ਗਿਆ} · \pa{ਗਈ} · \pa{ਗਏ} (giā · gaī · gae)}{ਗਿਆ · ਗਈ · ਗਏ (giā · gaī · gae)}]{Hear and say went again}

\label{lesson:PA-C216-gia-review}



\begin{quote}
\emph{Hear:} \emph{giā; gaī; gae} Say \textbf{\pa{ਗਿਆ} · \pa{ਗਈ} · \pa{ਗਏ}}.
\end{quote}

\subsection*{Your turn}
\begin{itemize}
\raggedright
  \item \emph{Hear:} \emph{the girl went home} Say \textbf{\pa{ਕੁੜੀ} \pa{ਘਰ} \pa{ਗਈ।}}
  \item \emph{Hear:} \emph{Aman went to school} Say \textbf{\pa{ਅਮਨ} \pa{ਸਕੂਲ} \pa{ਗਿਆ।}}
  \item \emph{Read it:} \textbf{\pa{ਗਏ}} Point to \textbf{\pa{ਏ}}, say \emph{e}, then say the whole form.
  \item \emph{Write it:} cover the model and write \textbf{\pa{ਗਈ}} once
\end{itemize}

\begin{tcolorbox}[breakable,colback=teal!4,colframe=teal!35!black,title={Before you move on}]
Which went has \textbf{\pa{ਏ}} at its end? (\textbf{\pa{ਗਏ}}, for a group.)
\end{tcolorbox}
\section[\texorpdfstring{\pa{ਕੱਲ੍ਹ} (kallh)}{ਕੱਲ੍ਹ (kallh)}]{\pa{ਕੱਲ੍ਹ} (kall) — yesterday; tomorrow}

\label{lesson:PA-C216-kall}



\begin{quote}
Before the new sense: say \emph{Aman went home}, then the farewell \textbf{\pa{ਕੱਲ੍ਹ} \pa{ਮਿਲਾਂਗੇ}}, \textquotedblleft{}we will meet tomorrow\textquotedblright{}.
\end{quote}

\subsection*{You'll want to know: \pa{ਕੱਲ੍ਹ}}
\textbf{\pa{ਕੱਲ੍ਹ}} — \emph{kall} — \textquotedblleft{}yesterday; tomorrow\textquotedblright{}. It names the day one step from today, either way; the verb tells you which.

\begin{quote}
\textbf{\pa{ਅਮਨ} \pa{ਕੱਲ੍ਹ} \pa{ਘਰ} \pa{ਗਿਆ।}} — \emph{Aman kallh ghar giā.} — Aman went home yesterday.
\end{quote}

With a went form it is yesterday; in \textbf{\pa{ਕੱਲ੍ਹ} \pa{ਮਿਲਾਂਗੇ}} it is tomorrow.

\begin{cousinweb}
\textbf{\pa{ਕੱਲ੍ਹ}} (\emph{kall}) is inherited from Sanskrit \emph{kālya}, \textquotedblleft{}at dawn; tomorrow\textquotedblright{}: CDIAL lists Punjabi \emph{kall}, \textquotedblleft{}yesterday, tomorrow\textquotedblright{} (CDIAL 3104).
\end{cousinweb}

\begin{grammarlens}[title={What you've built}]
One more word for talking about the past.
\end{grammarlens}

\subsection*{Your turn}
\begin{itemize}
\raggedright
  \item \emph{Say it:} \emph{kall}
  \item \emph{Hear:} \emph{the girl went to school yesterday} Say \textbf{\pa{ਕੁੜੀ} \pa{ਕੱਲ੍ਹ} \pa{ਸਕੂਲ} \pa{ਗਈ।}}
  \item \emph{Recall:} say \emph{we will meet tomorrow}, then \emph{Aman went home yesterday}
  \item \emph{Write it:} copy \textbf{\pa{ਕੱਲ੍ਹ}} once from the model
\end{itemize}

\begin{tcolorbox}[breakable,colback=teal!4,colframe=teal!35!black,title={Before you move on}]
What does \textbf{\pa{ਕੱਲ੍ਹ}} mean? (\textquotedblleft{}Yesterday; tomorrow\textquotedblright{}.) Which one, with \textbf{\pa{ਗਿਆ}}? (Yesterday.)
\end{tcolorbox}
\section[\texorpdfstring{\pa{ਉਹ} \pa{ਕੱਲ੍ਹ} \pa{ਕਿੱਥੇ} \pa{ਗਏ}? (uh kallh kitthe gae?)}{ਉਹ ਕੱਲ੍ਹ ਕਿੱਥੇ ਗਏ? (uh kallh kitthe gae?)}]{\pa{ਉਹ} \pa{ਕੱਲ੍ਹ} \pa{ਕਿੱਥੇ} \pa{ਗਏ}? — where did they go yesterday?}

\label{lesson:PA-C216-where-did-they-go}



\begin{quote}
Say \emph{yesterday}, then \emph{where?} Then \emph{the girl went to school}.
\end{quote}

\begin{grammarlens}[title={Grammar Lens: asking where}]
\textbf{\pa{ਕਿੱਥੇ}} goes right before the went form, where the place will go in the answer:

\begin{quote}
\textbf{\pa{ਉਹ} \pa{ਕੱਲ੍ਹ} \pa{ਕਿੱਥੇ} \pa{ਗਏ}?} — \emph{uh kallh kitthe gae?} — Where did they go yesterday?
\end{quote}

\begin{quote}
\textbf{\pa{ਉਹ} \pa{ਕੱਲ੍ਹ} \pa{ਬਾਜ਼ਾਰ} \pa{ਗਏ।}} — \emph{uh kallh bāzār gae.} — They went to the market yesterday.
\end{quote}

\begin{quote}
\textbf{\pa{ਅਮਨ} \pa{ਕੱਲ੍ਹ} \pa{ਕਿੱਥੇ} \pa{ਗਿਆ}?} — \emph{Aman kallh kitthe giā?} — Where did Aman go yesterday?
\end{quote}
\end{grammarlens}

\subsection*{Your turn}
\begin{itemize}
\raggedright
  \item \emph{Hear:} \emph{where did the girl go yesterday?} Say \textbf{\pa{ਕੁੜੀ} \pa{ਕੱਲ੍ਹ} \pa{ਕਿੱਥੇ} \pa{ਗਈ}?}
  \item \emph{Hear:} \emph{she went home} Say \textbf{\pa{ਉਹ} \pa{ਘਰ} \pa{ਗਈ।}}
  \item \emph{Read it:} \textbf{\pa{ਉਹ} \pa{ਕੱਲ੍ਹ} \pa{ਕਿੱਥੇ} \pa{ਗਏ}?} Point to \textbf{\pa{ਏ}} in \textbf{\pa{ਗਏ}}.
  \item \emph{Write it:} copy \textbf{\pa{ਉਹ} \pa{ਕੱਲ੍ਹ} \pa{ਕਿੱਥੇ} \pa{ਗਏ}?} once
\end{itemize}

\begin{tcolorbox}[breakable,colback=teal!4,colframe=teal!35!black,title={Before you move on}]
Ask where Aman went yesterday, then answer it.
\end{tcolorbox}
\section[\texorpdfstring{\pa{ਉਹ} \pa{ਕੱਲ੍ਹ} \pa{ਕਿੱਥੇ} \pa{ਗਏ}? (uh kallh kitthe gae?)}{ਉਹ ਕੱਲ੍ਹ ਕਿੱਥੇ ਗਏ? (uh kallh kitthe gae?)}]{Where they went: a four-skill payoff}

\label{lesson:PA-C216-payoff}



\begin{quote}
Say \emph{went} for a man, a woman and a group. Then \emph{yesterday}.
\end{quote}

\subsection*{Your turn: four skills}
\begin{itemize}
\raggedright
  \item \emph{Hear:} \emph{Aman went to school yesterday} Say \textbf{\pa{ਅਮਨ} \pa{ਕੱਲ੍ਹ} \pa{ਸਕੂਲ} \pa{ਗਿਆ।}}
  \item \emph{Hear:} \emph{where did they go?} Say \textbf{\pa{ਉਹ} \pa{ਕਿੱਥੇ} \pa{ਗਏ}?}
  \item \emph{Read it:} \textbf{\pa{ਕੁੜੀ} \pa{ਕੱਲ੍ਹ} \pa{ਘਰ} \pa{ਗਈ।}} Give its meaning. Then point to \textbf{\pa{ਏ}} in \textbf{\pa{ਗਏ}}.
  \item \emph{Write it:} cover the line and write \textbf{\pa{ਅਮਨ} \pa{ਕੱਲ੍ਹ} \pa{ਘਰ} \pa{ਗਿਆ।}}
\end{itemize}

\begin{tcolorbox}[breakable,colback=teal!4,colframe=teal!35!black,title={Before you move on}]
Say where you went yesterday, with the went form that fits you.
\end{tcolorbox}
